\problemname{Arborist}

Teo har äntligen landat sitt drömjobb; han har börjat jobba som arborist!
En viktig del av en arborists arbete är att plantera träd.

Teo behöver nu hjälp med att plantera $N$ träd på heltalspositioner längs en tallinje.
Tallinjen börjar vid position $0$ och värdet ökar med $1$ för varje steg åt höger (och minskar med $1$ åt vänster).

För alla $N$ träd har Teo redan bestämt exakt vilka heltalspositioner på tallinjen som de ska planteras på.
Dessa positioner är alla positiva heltal.

Teo befinner sig just nu vid position $0$, tillsammans med $N$ stora trädplantor.
För att plantera träden kan han göra följande:
\begin{itemize}
  \item Förflytta sig ett steg åt höger (öka positionen med $1$).
  \item Förflytta sig ett steg åt vänster (minska positionen med $1$).
  \item Om Teo befinner sig vid position $0$ kan han plocka upp en av de $N$ trädplantorna.
  \item Ta en av trädplantorna han bär på och plantera den vid positionen han befinner sig.
\end{itemize}
Varje sådan handling tar exakt en sekund.

Tyvärr klarar Teo av att bära som mest $K$ kilogram träd åt gången.
Varje träd väger antingen $20$ eller $40$ kilogram.
När alla träd har planterats måste Teo dessutom återvända till position $0$.

Din uppgift är att beräkna hur lång tid det tar för Teo att plantera alla träd och sedan återvända till position $0$, om han arbetar på snabbast möjliga sätt.

\section*{Indata}
Den första raden av indatan innehåller $N$ och $K$ ($1 \leq N \leq 8$, $20 \leq K \leq 50$),
antalet träd som ska planteras respektive antalet kilogram som Teo orkar bära.

Den andra raden innehåller $N$ heltal $x_1, \dots, x_N$ ($1 \leq x_i \leq 30$), positionerna
där varje trädplanta ska planteras.

Den tredje raden innehåller $N$ heltal $k_1, \dots, k_N$ ($k_i \in \{20, 40\}$, $k_i \leq K$),
vikten av varje trädplanta i kilogram.

Indatan beskriver alltså $N$ träd numrerade $1$ till $N$, där träd $i$ ska planteras på
position $x_i$ och väger $k_i$ kilogram.

\section*{Utdata}
Skriv ut ett heltal: minsta antalet sekunder Teo behöver för att plantera alla träd och
sedan återvända till position $0$.

\section*{Poängsättning}
Din lösning kommer att testas på en mängd testfallsgrupper.
För att få poäng för en grupp så måste du klara alla testfall i gruppen.

\noindent
\begin{tabular}{| l | l | p{12cm} |}
  \hline
  \textbf{Grupp} & \textbf{Poäng} & \textbf{Gränser} \\ \hline
  $1$    & $20$  & $K = 25$ \\ \hline
  $2$    & $20$  & $k_i = 20$ för alla $i$. \\ \hline
  $3$    & $20$  & Alla träd ska planteras på samma position. \\ \hline
  $4$    & $40$  & Inga ytterligare begränsningar. \\ \hline
\end{tabular}

\section*{Förklaring av exempelfall}
I exempelfall $1$ finns det två träd.
Ett av träden väger $20$ kilogram och ska planteras vid position $1$.
Det andra trädet väger $40$ kilogram och ska planteras vid position $2$.
En optimal lösning för Teo är att göra följande:
\begin{itemize}
  \item Plocka upp trädet som ska planteras vid position $1$.
  \item Gå höger.
  \item Plantera trädet som ska planteras vid position $1$.
  \item Gå vänster.
  \item Plocka upp trädet som ska planteras vid position $2$.
  \item Gå höger.
  \item Gå höger.
  \item Plantera trädet som ska planteras vid position $2$.
  \item Gå vänster.
  \item Gå vänster. Han har nu återvänt till position $0$ och är därmed färdig.
\end{itemize}
Det finns andra sätt att plantera träden lika snabbt, men det går inte att göra det på
mindre än $10$ sekunder.

I exempelfall $2$ är en optimal lösning att först plantera trädet vid position $9$, sedan
träden vid position $10$ och $8$ i en och samma tur, och till sist det vid position $7$.
